\documentclass[10pt,a4paper]{amsart}
\usepackage[margin=19mm]{geometry}
\usepackage{amsmath,amssymb,mathtools}
\usepackage[scr=rsfs,cal=euler]{mathalfa}
\usepackage{xfrac}
\usepackage[unicode,hidelinks,bookmarks=false]{hyperref}
\newcommand{\ie}{i.e.\ }
\newcommand{\cf}{cf.\ }
\newcommand{\resp}{respectively\ }
\newcommand{\Subsection}{\subsection}
\providecommand{\nobreakdash}{\nobreak\hbox{-}}
\setlength{\emergencystretch}{2em}
\multlinegap0pt
\usepackage{amssymb}
\usepackage{bm}
\usepackage{xcolor}
\newtheorem{lemma}{Lemma}
\def\di{\diamond}
\title{\bf Fundamental theorem of Poisson 3-Lie $(A,H)$-Hopf modules}
\author{Yan Ning, Daowei Lu, Dingguo Wang}
\date{Source: arXiv:2509.20670v2 (CC BY 4.0)}
\begin{document}
\maketitle

\noindent\emph{Source and licence.} \bf Fundamental theorem of Poisson 3-Lie $(A,H)$-Hopf modules. Version arXiv:2509.20670v2 (CC BY 4.0), distributed by arXiv under the Creative Commons Attribution 4.0 International licence, http://creativecommons.org/licenses/by/4.0/, as recorded in the arXiv licence metadata for this exact version and shown on the abstract page https://arxiv.org/abs/2509.20670v2. Full licence notice: \texttt{licenses/arxiv-2509.20670-CC-BY-4.0.txt}.\par\smallskip

\noindent\textit{Editorial scope.} Counted unit: the article's lemma lem:3d (source0.tex lines 515--524). The article's complete original proof of it is delivered with it (source0.tex lines 526--578, one \texttt{proof} environment of 53 lines). No mathematics is written, repaired or re-derived: the statement and the proof are the article's own lines, in the article's own notation, and no retained line is substituted or re-flowed.  Retained uncounted as the source-local context the proof needs: lines 156--160.  Nothing was omitted from any retained span, so the package carries no omission record.  No retained line contains a citation, so the wrapper carries no bibliography and no bibliography entry.  No retained line contains a figure environment, an image-inclusion command or drawing code, so no image is delivered and the package declares no asset dependency.  Editorial formatting only, in addition: an AMS \texttt{amsart} preamble that restates the article's own declarations and macros used by the retained lines, namely the environments \texttt{lemma} on separate counters, together with the standard packages those lines need.  The retained environments print in this document as Lemma 1 in order of appearance, whereas the article prints the same items with the numbering of its own full document.  The complete native source of the article as distributed in the arXiv e-print for this exact version is bundled unchanged as \texttt{sources/185693/source0.tex}.
\begin{align}
&(x, y)\di m=-(y, x)\di m,\label{1b}\\
&(x, y)\di((u,v)\di m)-(u,v)\di((x, y)\di m)=(\{x,y,u\},v)\di m+(u,\{x,y,v\})\di m,\label{1c}\\
&(\{x,y,u\},v)\di m=(y,u)\di((x,v)\di m)-(x,u)\di((y,v)\di m)+(x,y)\di((u,v)\di m),\label{1d}
\end{align}

\begin{lemma}\label{lem:3d}
Let $H$ be a weak Hopf algebra. Let $M\in\! _{\mathcal{P}A}\mathcal{M}^H$. Suppose that there exists a right $H$-colinear algebra map $\phi:H\rightarrow A^A$, then for all $x,y,a\in A,m\in M$,
\begin{itemize}
  \item [(1)] $p_M(a\cdot m)=p_A(a)\cdot p_M(m),$
  \item [(2)] $(x,y)\di'p_M(m)=(x,y)\di' m,$
  \item [(3)] $(M,\di')$ is a 3-Lie $A$-module,
  \item [(4)] $(x,y)\di p_M(m)=\phi(x_{(1)}y_{(1)})\cdot ((x_{(0)},y_{(0)})\di'p_M(m))$,
  \item [(5)] $\phi(m_{(1)})\cdot p_M(m_{(0)})=m$.
\end{itemize}
\end{lemma}

\begin{proof}
We only verify the conditions (2),(3) and (4).

(2) For all $x,y\in A,m\in M$,
\begin{align*}
(x,y)\di'p_M(m)&=p_M((x,y)\di p_M(m))\\
&=p_M((x,y)\di (\phi(S^{-1}(m_{(1)}))\cdot m_{(0)}))\\
&=p_M(\{x,y,\phi(S^{-1}(m_{(1)}))\}\cdot m_{(0)}+\phi(S^{-1}(m_{(1)}))\cdot((x,y)\di m_{(0)}))\\
&=p_M(\phi(S^{-1}(m_{(1)}))\cdot((x,y)\di m_{(0)}))\\
&=p_A(\phi(S^{-1}(m_{(1)})))\cdot p_M((x,y)\di m_{(0)})\\
%&=\phi(S^{-1}(\phi(S^{-1}(m_{(1)})_2)))\phi(S^{-1}(m_{(1)})_1)\cdot p_M((x,y)\di m_{(0)})\\
&=\phi(S^{-2}(m_{(1)2}))\phi(S^{-1}(m_{(1)3}))\cdot p_M((x,y)\di m_{(0)})\\
&=p_M((x,y)\di m)\\
&=(x,y)\di' m.
\end{align*}

(3) For all $x,y,u,v\in A,m\in M$, 
$$(x,y)\di' m=p_M((x,y)\di m)=-p_M((y,x)\di m)=-(y,x)\di' m.$$
Then
\begin{align*}
&(\{x,y,u\},v)\di' m+(u,\{x,y,v\})\di' m\\
&=(\{x,y,u\},v)\di' p_M(m)+(u,\{x,y,v\})\di' p_M(m)\\
&=p_M((\{x,y,u\},v)\di p_M(m))+p_M((u,\{x,y,v\})\di p_M(m))\\
&\stackrel{(\ref{1c})}{=}p_M((x, y)\di((u,v)\di p_M(m))-(u,v)\di((x, y)\di p_M(m)))\\
&=(x, y)\di'((u,v)\di p_M(m))-(u,v)\di'((x, y)\di p_M(m))\\
&=(x, y)\di' p_M((u,v)\di p_M(m))-(u,v)\di'p_M((x, y)\di p_M(m))\\
&=(x, y)\di' ((u,v)\di' p_M(m))-(u,v)\di'((x, y)\di' p_M(m))\\
&=(x, y)\di'((u,v)\di' m)-(u,v)\di'((x, y)\di' m),
\end{align*}
and
\begin{align*}
&(\{x,y,u\},v)\di' m=(\{x,y,u\},v)\di' p_M(m)=p_M((\{x,y,u\},v)\di p_M(m))\\
&\stackrel{(\ref{1d})}{=}p_M((y,u)\di((x,v)\di p_M(m)))-p_M((x,u)\di((y,v)\di p_M(m)))+p_M((x,y)\di((u,v)\di p_M(m)))\\
&=(y,u)\di'((x,v)\di p_M(m))-(x,u)\di'((y,v)\di p_M(m))+(x,y)\di'((u,v)\di p_M(m))\\
&=(y,u)\di'p_M((x,v)\di p_M(m))-(x,u)\di'p_M((y,v)\di p_M(m))+(x,y)\di'p_M((u,v)\di p_M(m))\\
&=(y,u)\di'((x,v)\di' p_M(m))-(x,u)\di'((y,v)\di' p_M(m))+(x,y)\di'((u,v)\di' p_M(m))\\
&=(y,u)\di'((x,v)\di' m)-(x,u)\di'((y,v)\di' m)+(x,y)\di'((u,v)\di' m).
\end{align*}
Therefore $(M,\di')$ is a 3-Lie $A$-module.

(4) Let $x,y\in A,m\in M$,
\begin{align*}
&\phi(x_{(1)}y_{(1)})\cdot ((x_{(0)},y_{(0)})\di'p_M(m))\\
&=\phi(x_{(1)}y_{(1)})\cdot ((x_{(0)},y_{(0)})\di'm)\\
&=\phi(x_{(1)}y_{(1)})\cdot p_M(((x_{(0)},y_{(0)})\di m))\\
&=\phi(x_{(1)}y_{(1)})\phi(S^{-1}(((x_{(0)},y_{(0)})\di m)_{(1)}))\cdot (x_{(0)},y_{(0)})\di m)_{(0)}\\
&=\phi(x_{(1)2}y_{(1)2})\phi(S^{-1}(x_{(1)1}y_{(1)1}m_{(1)}))\cdot ((x_{(0)},y_{(0)})\di m_{(0)})\\
&=\phi(S^{-1}(m_{(1)}))\cdot ((x,y)\di m_{(0)})\\
&=(x,y)\di(\phi(S^{-1}(m_{(1)}))\cdot m_{(0)})-\{x,y,\phi(S^{-1}(m_{(1)}))\}\cdot m_{(0)}\\
&=(x,y)\di p_M(m).
\end{align*}
The proof is completed.
\end{proof}

\end{document}
